% Standalone problem document, generated from the adjacent README.md.
% Compile with XeLaTeX. Mathematical content is maintained in Markdown.
\documentclass[11pt,a4paper]{article}
\usepackage[fontsize=10.5pt]{scrextend}
\usepackage[margin=22mm,headheight=15pt,headsep=9mm,footskip=11mm]{geometry}
\usepackage{fontspec}
\setmainfont{texgyrepagella}[Extension=.otf,UprightFont=*-regular,BoldFont=*-bold,ItalicFont=*-italic,BoldItalicFont=*-bolditalic]
\setsansfont{texgyreheros}[Extension=.otf,UprightFont=*-regular,BoldFont=*-bold,ItalicFont=*-italic,BoldItalicFont=*-bolditalic]
\setmonofont{texgyrecursor}[Extension=.otf,UprightFont=*-regular,BoldFont=*-bold,ItalicFont=*-italic,BoldItalicFont=*-bolditalic,Scale=MatchLowercase]
\usepackage{amsmath,amssymb}
\usepackage{unicode-math}
\setmathfont{texgyrepagella-math.otf}
\usepackage{xcolor,microtype,enumitem,needspace,fancyhdr}
\usepackage{bookmark}
\definecolor{ink}{HTML}{17344B}
\definecolor{accent}{HTML}{197D80}
\definecolor{muted}{HTML}{596773}
\hypersetup{colorlinks=true,linkcolor=accent,urlcolor=accent,pdftitle={MI-20 - Sharp
subquadratic Lee constants for sums of
matrices},pdfauthor={Open Problems in Numerical Linear Algebra}}
\urlstyle{same}
\setlength{\parindent}{0pt}
\setlength{\parskip}{5pt plus 1pt minus 1pt}
\linespread{1.04}
\setlength{\emergencystretch}{3em}
\widowpenalty=10000
\clubpenalty=10000
\displaywidowpenalty=10000
\allowdisplaybreaks[1]
\setlist{nosep,leftmargin=1.5em}
\providecommand{\tightlist}{\setlength{\itemsep}{0pt}\setlength{\parskip}{0pt}}
\providecommand{\pandocbounded}[1]{#1}
\pagestyle{fancy}
\fancyhf{}
\fancyhead[L]{\sffamily\footnotesize\color{muted}OPEN PROBLEMS IN NUMERICAL LINEAR ALGEBRA}
\fancyhead[R]{\sffamily\bfseries\footnotesize\color{accent}MI-20}
\fancyfoot[L]{\parbox[b]{0.9\textwidth}{\sffamily\fontsize{8}{10}\selectfont\color{muted}Editorial ratings; a bounded search cannot certify that a problem remains unsolved.\\Literature check: 2026-09-12}}
\fancyfoot[R]{\sffamily\footnotesize\color{muted}\thepage}
\renewcommand{\headrulewidth}{0.3pt}
\renewcommand{\headrule}{\hbox to\headwidth{\color{accent}\leaders\hrule height \headrulewidth\hfill}}
\makeatletter
\renewcommand{\section}{\Needspace{5\baselineskip}\@startsection{section}{1}{0pt}{12pt plus 2pt minus 2pt}{4pt}{\sffamily\bfseries\large\color{ink}}}
\renewcommand{\subsection}{\Needspace{4\baselineskip}\@startsection{subsection}{2}{0pt}{9pt plus 1pt minus 1pt}{3pt}{\sffamily\bfseries\normalsize\color{ink}}}
\makeatother
\setcounter{secnumdepth}{0}
\begin{document}
{\sffamily\small\color{muted}Matrix inequalities and norms\par}
\vspace{3pt}
{\raggedright\hyphenpenalty=10000\exhyphenpenalty=10000\sffamily\bfseries\fontsize{21}{24}\selectfont\color{ink}Sharp
subquadratic Lee constants for sums of matrices\par}
\vspace{7pt}
{\sffamily\small\textbf{Difficulty:} challenging\quad\textbf{Importance:} interesting
to the community\par}
{\sffamily\small\bfseries\color{accent}Status: Open\par}
\vspace{4pt}
{\color{accent}\hrule height 0.8pt}
\vspace{6pt}
\textbf{Rating rationale:} Determining a sharp two-parameter function
after the proposed formula failed is challenging; matrix-sum bounds
affect norm estimates used across NLA.

\subsection{Audited dual and projective reductions --- 12 September
2026}\label{audited-dual-and-projective-reductions-12-september-2026}

\textbf{Author:} Sidney Holden, Center for Computational Biology,
Flatiron Institute, Simons Foundation.
\href{https://github.com/ajt60gaibb/OpenProblemsInNLA/blob/main/references/holden-matrix-2026-09-12/README.md}{Submission
and verified affiliation}.

\href{https://github.com/ajt60gaibb/OpenProblemsInNLA/blob/main/references/holden-matrix-2026-09-12/MI-20/result.md}{Sections
1--4} give the known positive dual formulation (Qiu, Proposition 3.1), a
fixed-dimension binary projection reduction, and an all-summands
projective dilation preserving the supremum over dimensions. These
lemmas do not determine \(C_p(m)\) for the requested parameter range;
the target remains Open.

A separate
\href{https://github.com/ajt60gaibb/OpenProblemsInNLA/blob/main/references/holden-matrix-2026-09-12/verification/MI-20-MI-27-review.md}{independent
Codex AI-agent audit} passed this limited scope. This is informal
automated review; no complete resolution, historical novelty, external
human peer review or formal verification is claimed. No Lean
verification was performed.

\subsection{Problem statement}\label{problem-statement}

For a complex square matrix \(X\), write \(|X|=(X^*X)^{1/2}\) and
\(\|X\|_p=(\mathop{\mathrm{tr}}\nolimits|X|^p)^{1/p}\). For every
integer \(m\ge2\) and real \(1< p<2\), determine exactly the
dimension-independent constant

\[C_p(m)=\sup_{n\ge1}\ \sup_{A_1,\ldots,A_m\in\mathbb C^{n\times n},\,\sum_j|A_j|\ne0}
\frac{\|\sum_{j=1}^m A_j\|_p}{\|\sum_{j=1}^m|A_j|\|_p}.\]

Thus the requested answer is the optimal constant, as a function of both
\(p\) and \(m\), in the corresponding Schatten norm inequality. An upper
bound alone does not settle the problem. All matrix orders are included
in the same supremum.

These constants quantify the loss when replacing a matrix sum by the sum
of its positive semidefinite moduli, an operation common in norm
estimates for matrix computations. This is distinct from constants
comparing the arithmetic and symmetric moduli of one matrix (catalog
MI-09).

\subsection{References}\label{references}

\begin{enumerate}
\def\labelenumi{\arabic{enumi}.}
\tightlist
\item
  X. Li, \emph{Sharp Concave-Function Transfer for Lee-Type Schatten
  Norm Inequalities}, arXiv:2608.25989v1 (2026), §1, equations defining
  the linear constants, and §4, Problem 4.2.
  \href{https://arxiv.org/html/2608.25989}{Full text}.
\item
  Q. Tang and S. Zhang, \emph{Generalizing Lee's conjecture on the sum
  of absolute values of matrices}, Linear Algebra Appl. 731 (2026),
  196--204, original candidate and endpoint results.
  \href{https://arxiv.org/abs/2510.16846}{Preprint}.
\item
  H. Qiu, \emph{Sharp quasi-reverse Minkowski inequality for Schatten
  norms} (2026), the two-matrix sharp result for \(p\ge2\) and
  counterexamples in the subquadratic interval.
  \href{https://arxiv.org/abs/2608.17565}{Preprint}.
\end{enumerate}

Status check (2026-09-10): Li's current v1, dated August 26, explicitly
retains this problem after proving equality of the sharp linear and
concave-function constants. Its §4 records failure of the Tang--Zhang
candidate throughout \(1< p<2\). Searches for the exact title, ``Lee
sharp Schatten constants'', and ``Tang Zhang conjecture'' with 2025/2026
found no later determination of this entire function. This is a bounded
literature check, not a guarantee against an unindexed solution.
\end{document}
